\documentclass[11pt]{article}
% Shared preamble for all solution files.
\usepackage[margin=1.1in]{geometry}
\usepackage{enumitem}
\usepackage{tikz}
\usepackage[hidelinks]{hyperref}
\usepackage{amsmath,amssymb,amsthm}
\newcommand{\fa}{\mathfrak a}
\newcommand{\fb}{\mathfrak b}
\newcommand{\fc}{\mathfrak c}
\newcommand{\fp}{\mathfrak p}
\newcommand{\fq}{\mathfrak q}
\newcommand{\fm}{\mathfrak m}
\newcommand{\fn}{\mathfrak n}
\newcommand{\fr}{\mathfrak r}
\newcommand{\Nil}{\mathfrak N}
\newcommand{\Jac}{\mathfrak R}
\newcommand{\rad}{r}
\newcommand{\Spec}{\operatorname{Spec}}
\newcommand{\Max}{\operatorname{Max}}
\newcommand{\Ann}{\operatorname{Ann}}
\newcommand{\Ker}{\operatorname{Ker}}
\newcommand{\Z}{\mathbf Z}
\newcommand{\Q}{\mathbf Q}
\newcommand{\R}{\mathbf R}
\newcommand{\C}{\mathbf C}
\newcommand{\ci}{\mathrm{c}}
\newcommand{\ex}{\mathrm{e}}

\setlength{\parindent}{0pt}
\setlength{\parskip}{0.55em}
% \exercise{label}{statement}  -- followed by the solution text
\newcommand{\exercise}[2]{\par\bigskip\noindent{\large\textbf{Exercise #1}}\par\smallskip
  \begin{quote}\small\textbf{Problem.}\ #2\end{quote}\par}
\newcommand{\pf}{\noindent\textbf{Solution.}\ }

\title{Atiyah--Macdonald, Chapter 9: Discrete Valuation Rings and Dedekind Domains\\[2pt]\large Solutions to the exercises}
\author{}
\date{}
\begin{document}
\maketitle
\noindent\textbf{Conventions.} $K$ denotes the field of fractions of the domain $A$. A \emph{Dedekind domain} is a Noetherian integrally closed domain of dimension one (Thm.~9.3); in particular it is not a field, every non-zero prime is maximal, and $A_\fp$ is a discrete valuation ring (DVR) for every non-zero prime $\fp$. From the text I use: 9.2 (a Noetherian local domain of dimension one with principal maximal ideal is a DVR; in a DVR with maximal ideal $(\pi)$ the non-zero ideals are exactly $(\pi^k)$, $k\ge0$); 9.3, 9.4 (unique factorization of non-zero ideals into primes); 9.6--9.9 (a non-zero fractional ideal of a Dedekind domain is invertible; the non-zero fractional ideals form a group $I$, free on the non-zero primes; $H=I/P$ is the ideal class group); the localization facts of Ch.~3 (3.3, 3.4, 3.8, 3.9, 3.10, 3.11); integral closure commutes with localization (5.12); the Chinese remainder theorem (1.10, 1.16); and the structure theorem for finitely generated modules over a principal ideal domain, including uniqueness of the elementary divisors.

A recurring principle (3.8/3.9): \emph{if $X\subseteq Y$ are ideals (or modules) with $X_\fm=Y_\fm$ for every maximal ideal $\fm$, then $X=Y$} (apply 3.8 to $Y/X$, using exactness of localization). In a Dedekind domain the local rings $A_\fm$ at maximal ideals are DVRs, so identities can be verified in a DVR.

This chapter has no numbered in-text exercises (the text only says ``check that $v$ is well defined'' in the proof of 9.2), and the 9 end-of-chapter exercises are all solved below.

%=====================================================================
\section*{End-of-chapter exercises}

%---------------------------------------------------------------- 1
\exercise{1}{Let $A$ be a Dedekind domain, $S$ a multiplicatively closed subset of $A$. Show that $S^{-1}A$ is either a Dedekind domain or the field of fractions of $A$.

Suppose that $S\ne A-\{0\}$, and let $H,H'$ be the ideal class groups of $A$ and $S^{-1}A$ respectively. Show that extension of ideals induces a surjective homomorphism $H\to H'$.}
\pf
We assume $0\notin S$ (otherwise $S^{-1}A=0$ is not a domain). Then $A\subseteq S^{-1}A\subseteq K$ and $S^{-1}A$ is a domain with field of fractions $K$.

\emph{Noetherian.} Every ideal $\fb$ of $S^{-1}A$ is extended, $\fb=(\fb^c)^e$ (3.11(i)); if $\fb^c$ is generated by $x_1,\dots,x_n$, then $\fb$ is generated by $x_1/1,\dots,x_n/1$.

\emph{Integrally closed.} $A$ is integrally closed, so $S^{-1}A$ is integrally closed in $K$ (5.12).

\emph{Primes.} The primes of $S^{-1}A$ correspond bijectively, preserving inclusion, to the primes of $A$ not meeting $S$ (3.11(iv)). The prime $(0)$ always corresponds to $(0)$ ($0\cap S=\emptyset$), and every non-zero prime of $A$ is maximal. Hence every chain of primes of $S^{-1}A$ has length at most $1$.

\emph{Two cases.} If every non-zero prime of $A$ meets $S$, then $(0)$ is the only prime of $S^{-1}A$; a non-unit of $S^{-1}A$ lies in some maximal ideal, which is a non-zero prime unless the element is $0$; so every non-zero element is a unit and $S^{-1}A=K$. Otherwise $S^{-1}A$ has a non-zero prime, so $\dim S^{-1}A=1$, and $S^{-1}A$ is a Noetherian, integrally closed domain of dimension one, i.e.\ a Dedekind domain.

\emph{The map $H\to H'$.} (If $S^{-1}A=K$ the group $H'$ is trivial, since the only non-zero fractional ideal of $K$ is $K$, and there is nothing to prove; so assume $S^{-1}A$ is Dedekind.) Let $I,I'$ be the groups of non-zero fractional ideals of $A$ and $S^{-1}A$. For $M\in I$ put $\theta(M)=S^{-1}M=S^{-1}A\cdot M\subseteq K$. If $xM\subseteq A$ with $0\ne x\in A$ then $x\,S^{-1}M\subseteq S^{-1}A$, so $\theta(M)$ is a non-zero fractional ideal. Clearly $\theta(MN)=\theta(M)\theta(N)$ (products are generated by the products of generators), so $\theta:I\to I'$ is a group homomorphism (it preserves identities and inverses because $\theta(A)=S^{-1}A$). For $u\in K^*$, $\theta((u))=(u)$, so $\theta$ maps the principal fractional ideals $P$ into $P'$ and induces $H=I/P\to H'=I'/P'$.

\emph{Surjective.} It suffices to show $\theta:I\to I'$ is surjective. Let $N'\in I'$, and choose $0\ne z\in S^{-1}A$ with $zN'\subseteq S^{-1}A$; writing $z=a/s$ we may replace $z$ by $a=sz\in A-\{0\}$ ($s$ is a unit of $S^{-1}A$). Then $\fb=aN'$ is a non-zero ideal of $S^{-1}A$, hence $\fb=\fa^e=S^{-1}\fa$ with $\fa=\fb^c$ an ideal of $A$ (3.11(i)). Therefore $N'=a^{-1}S^{-1}\fa=S^{-1}(a^{-1}\fa)$, and $M=a^{-1}\fa$ is a non-zero fractional ideal of $A$ ($aM=\fa\subseteq A$). So $N'=\theta(M)$. \hfill$\square$

%---------------------------------------------------------------- 2
\exercise{2}{Let $A$ be a Dedekind domain. If $f=a_0+a_1x+\cdots+a_nx^n$ is a polynomial with coefficients in $A$, the \emph{content} of $f$ is the ideal $c(f)=(a_0,\dots,a_n)$ in $A$. Prove \emph{Gauss's lemma} that $c(fg)=c(f)c(g)$. [Localize at each maximal ideal.]}
\pf
Always $c(fg)\subseteq c(f)c(g)$, because every coefficient of $fg$ is a sum of products $a_ib_j$ of coefficients of $f$ and $g$. By the principle in the Conventions it suffices to show $c(fg)_\fm=(c(f)c(g))_\fm$ for every maximal ideal $\fm$. Since localization commutes with products (3.4), $(c(f)c(g))_\fm=c(f)_\fm c(g)_\fm$, and $c(f)_\fm$ is the ideal of $A_\fm$ generated by the (images of the) coefficients of $f$, i.e.\ the content of $f$ viewed in $A_\fm[x]$; likewise for $g$ and $fg$. So we may replace $A$ by $A_\fm$ and assume $A$ is a DVR with maximal ideal $(\pi)$ and valuation $v$ (recall $\fm\ne0$ since $A$ is not a field).

If $f=0$ or $g=0$ both sides are $0$. Otherwise let $a=\min_iv(a_i)$ and $b=\min_jv(b_j)$. Then $c(f)=(\pi^a)$, $c(g)=(\pi^b)$ (ideals of a DVR are the $(\pi^k)$ and $\pi^a$ divides every coefficient). Write $f=\pi^af_1$, $g=\pi^bg_1$ with $f_1,g_1\in A[x]$; then some coefficient of $f_1$ and some coefficient of $g_1$ is a unit, so $\bar f_1\ne0$ and $\bar g_1\ne0$ in $(A/\fm)[x]$. As $(A/\fm)[x]$ is a domain, $\bar f_1\bar g_1\ne0$, i.e.\ some coefficient of $f_1g_1$ is not in $\fm$, so is a unit, and $c(f_1g_1)=(1)$. Hence $c(fg)=\pi^{a+b}c(f_1g_1)=(\pi^{a+b})=c(f)c(g)$. \hfill$\square$

%---------------------------------------------------------------- 3
\exercise{3}{A valuation ring (other than a field) is Noetherian if and only if it is a discrete valuation ring.}
\pf
Recall that a \emph{valuation ring} is an integral domain $A$ with field of fractions $K$ such that for every $x\in K^*$ either $x\in A$ or $x^{-1}\in A$.

($\Leftarrow$) A DVR is Noetherian: its ideals are $0$ and $\fm^k=(x^k)$, all finitely generated (this is shown in the text before 9.2).

($\Rightarrow$) Let $A$ be a Noetherian valuation ring which is not a field.

\emph{$A$ is local.} Let $x,y$ be non-units of $A$, not both $0$, say $y\ne0$. Then $x/y\in A$ or $y/x\in A$; by symmetry (if $x\ne0$) suppose $a=x/y\in A$. Then $x+y=y(1+a)$ is a multiple of the non-unit $y$, hence a non-unit. Also $ax$ is a non-unit for $a\in A$, $x$ a non-unit. So the non-units form an ideal, which is therefore the unique maximal ideal $\fm$, and $\fm\ne0$ because $A$ is not a field.

\emph{Finitely generated ideals are principal.} For $a,b\in A$ not both $0$ say $b\ne0$ and $a/b\in A$ (else $b/a\in A$ and swap the roles); then $a\in(b)$ and $(a,b)=(b)$. By induction every finitely generated ideal of $A$ is principal. $A$ is Noetherian, so $A$ is a principal ideal domain.

\emph{Dimension one.} A non-zero prime ideal in a PID is maximal, so the only non-zero prime is $\fm$, and $\dim A=1$.

Thus $A$ is a Noetherian local domain of dimension one whose maximal ideal is principal, so $A$ is a DVR by 9.2 (iii)$\Rightarrow$(i). \hfill$\square$

%---------------------------------------------------------------- 4
\exercise{4}{Let $A$ be a local domain which is not a field and in which the maximal ideal $\fm$ is principal and $\bigcap_{n=1}^\infty\fm^n=0$. Prove that $A$ is a discrete valuation ring.}
\pf
Let $\fm=(x)$; $x\ne0$ since $A$ is not a field. Let $a\in A$, $a\ne0$. As $\bigcap\fm^n=0$ there is a largest $k\ge0$ with $a\in\fm^k=(x^k)$; write $a=x^ku$. Then $u\notin\fm$ (else $a\in(x^{k+1})$), so $u$ is a unit. The integer $k=:v(a)$ is uniquely determined by $a$: if $x^ku=x^lw$ with $k<l$ and $u,w$ units, cancelling $x^k$ in the domain gives $u=x^{l-k}w\in\fm$, absurd.

Every element of $K^*$ is $a/b$ with $a,b\in A-\{0\}$; define $v(a/b)=v(a)-v(b)$. This is well defined: if $a/b=a'/b'$ then $ab'=a'b$, so $v(a)+v(b')=v(a')+v(b)$ (units do not change $v$, and $v(ab)=v(a)+v(b)$ on $A-\{0\}$ since $x^ku\cdot x^lw=x^{k+l}uw$). Thus $v:K^*\to\mathbf Z$ is a homomorphism, and it is onto because $v(x)=1$.

\emph{Ultrametric inequality.} Let $\alpha=a/c$, $\beta=b/c$ with common denominator $c$ (put $v(0)=+\infty$; if $\alpha+\beta=0$, or if $\alpha$ or $\beta$ is $0$, there is nothing to prove). Say $v(a)\le v(b)$ and $k=v(a)$. Then $a,b\in(x^k)$, so $a+b\in(x^k)$ and $v(a+b)\ge k$ (if $a+b\ne0$). Hence $v(\alpha+\beta)=v(a+b)-v(c)\ge v(a)-v(c)=\min(v(\alpha),v(\beta))$. So $v$ is a discrete valuation.

\emph{$A$ is its valuation ring.} If $z=a/b$ with $v(z)\ge0$, i.e.\ $v(a)\ge v(b)$, then $a=x^ku$, $b=x^lw$ with $k\ge l$, so $z=x^{k-l}uw^{-1}\in A$. Conversely for $a\in A-\{0\}$, $v(a)\ge0$. Hence $A=\{z:v(z)\ge0\}\cup$ ($0$) is the valuation ring of $v$, and $A$ is a discrete valuation ring by definition. \hfill$\square$

%---------------------------------------------------------------- 5
\exercise{5}{Let $M$ be a finitely generated module over a Dedekind domain. Prove that $M$ is flat $\iff$ $M$ is torsion-free. [Use Chapter 3, Exercise 13 and Chapter 7, Exercise 16.]}
\pf
Let $A$ be the Dedekind domain. (I give a self-contained argument that does not use Chapter 7, Exercise 16; I only use the structure theorem for finitely generated modules over a PID.)

($\Rightarrow$) Let $M$ be flat and $0\ne a\in A$. Multiplication by $a$, $A\to A$, is injective; tensoring with the flat module $M$ gives that $M\to M$, $m\mapsto am$, is injective. So no non-zero element of $M$ is killed by a non-zero element of $A$: $M$ is torsion-free. (This direction holds for any domain and any module.)

($\Leftarrow$) Let $M$ be finitely generated and torsion-free. By Ch.~3, Ex.~13, $M_\fm$ is torsion-free for every maximal ideal $\fm$. Here $\fm\ne0$ and $A_\fm$ is a DVR, hence a PID. $M_\fm$ is a finitely generated torsion-free module over the PID $A_\fm$, hence free (structure theorem: $M_\fm\cong A_\fm^r\oplus T$ with $T$ the torsion submodule, and $T=0$). A free module is flat. Thus $M_\fm$ is a flat $A_\fm$-module for all maximal ideals $\fm$, and so $M$ is flat over $A$ by 3.10. \hfill$\square$

%---------------------------------------------------------------- 6
\exercise{6}{Let $M$ be a finitely generated torsion module ($T(M)=M$) over a Dedekind domain $A$. Prove that $M$ is uniquely representable as a finite direct sum of modules $A/\fp_i^{n_i}$, where $\fp_i$ are non-zero prime ideals of $A$. [For each $\fp\ne0$, $M_\fp$ is a torsion $A_\fp$-module; use the structure theorem for modules over a principal ideal domain.]}
\pf
``Uniquely'' means: the multiset of pairs $(\fp_i,n_i)$ (with $n_i\ge1$) is determined by $M$.

\emph{$\Ann(M)\ne0$.} If $x_1,\dots,x_n$ generate $M$ and $0\ne a_i\in\Ann(x_i)$, then $a=a_1\cdots a_n\ne0$ kills $M$. Let $\fa=\Ann(M)\ne0$. If $\fa=A$ then $M=0$ (the empty sum); otherwise by 9.4 $\fa=\fp_1^{e_1}\cdots\fp_r^{e_r}$ with distinct non-zero primes $\fp_i$ and $e_i\ge1$. Put $\fq_i=\fp_i^{e_i}$ and $\fb_i=\prod_{j\ne i}\fq_j$.

\emph{Existence.} The ideals $\fq_i$ are pairwise coprime (distinct maximal ideals have coprime powers: a maximal ideal containing $\fq_i+\fq_j$ would contain $\fp_i$ and $\fp_j$, forcing $\fp_i=\fp_j$), so $\fb_1+\dots+\fb_r=(1)$ (a maximal ideal containing all $\fb_i$ would, for $r\ge2$, contain some $\fq_j$, hence equal some $\fp_j$ with $j\ne1$, and also some $\fq_k$, $k\ne j$, hence equal $\fp_k\ne\fp_j$; for $r=1$, $\fb_1=A$). Write $1=\sum c_ib_i$ with $b_i\in\fb_i$, and put $M_i=\{x\in M:\fq_ix=0\}$. For $x\in M$, $x=\sum c_ib_ix$ and $c_ib_ix\in M_i$ because $\fq_ib_i\subseteq\fa$. The sum is direct: if $\sum y_i=0$ with $y_i\in M_i$, multiply by $b_i$: $b_iy_j=0$ for $j\ne i$ (since $\fq_j\mid\fb_i$ there) so $b_iy_i=0$; as also $\fq_iy_i=0$ and $\fb_i+\fq_i=A$ ($\fq_i$ is coprime to $\fb_i$), $y_i=0$. So $M=\bigoplus_{i=1}^rM_i$, each $M_i$ finitely generated (a quotient of $M$) and killed by $\fq_i=\fp_i^{e_i}$.

Fix $i$ and write $\fp=\fp_i$, $\fq=\fq_i$. For $s\in A-\fp$ we have $(s)+\fq=A$ (the only maximal ideal containing $\fq$ is $\fp$), so $ts\equiv1\pmod\fq$ for some $t$, and $s$ acts bijectively on $M_i$. Hence $M_i\to(M_i)_\fp$, $x\mapsto x/1$, is an isomorphism of $A$-modules (injective since $sx=0\Rightarrow x=tsx=0$; surjective since $x/s=tx/1$). Now $(M_i)_\fp$ is a finitely generated module over the DVR $A_\fp$, killed by $\fp^{e}A_\fp$, so it is a torsion module and by the structure theorem
\[
(M_i)_\fp\cong\bigoplus_jA_\fp/\fp^{k_j}A_\fp\qquad(k_j\ge1),
\]
and $A_\fp/\fp^kA_\fp\cong(A/\fp^k)_\fp\cong A/\fp^k$ (by the same argument as above, elements outside $\fp$ act invertibly on $A/\fp^k$). Therefore $M_i\cong\bigoplus_jA/\fp_i^{k_{ij}}$ as $A$-modules, and $M\cong\bigoplus_{i,j}A/\fp_i^{k_{ij}}$.

\emph{Uniqueness.} Suppose $M\cong\bigoplus_{j=1}^mA/\fp_j^{n_j}$ with non-zero primes $\fp_j$ and $n_j\ge1$. Fix a non-zero prime $\fp$. If $\fp_j\ne\fp$, then $\fp_j^{n_j}\not\subseteq\fp$ (else $\fp_j\subseteq\fp$, so $\fp_j=\fp$ as both are maximal), so there is $s\in\fp_j^{n_j}-\fp$, which kills $A/\fp_j^{n_j}$, and hence $(A/\fp_j^{n_j})_\fp=0$. If $\fp_j=\fp$ then $(A/\fp^{n_j})_\fp\cong A_\fp/\fp^{n_j}A_\fp$ (3.4). Thus
\[
M_\fp\cong\bigoplus_{j:\ \fp_j=\fp}A_\fp/\fp^{n_j}A_\fp .
\]
By the uniqueness part of the structure theorem over the PID $A_\fp$ (the elementary divisors $\fp^{n_j}A_\fp$, equivalently the exponents $n_j$, are determined by $M_\fp$), the multiset $\{n_j:\fp_j=\fp\}$ depends only on $M$. This holds for every non-zero prime $\fp$, so the multiset of pairs $(\fp_j,n_j)$ is determined by $M$. \hfill$\square$

%---------------------------------------------------------------- 7
\exercise{7}{Let $A$ be a Dedekind domain and $\fa\ne0$ an ideal in $A$. Show that every ideal in $A/\fa$ is principal.

Deduce that every ideal in $A$ can be generated by at most $2$ elements.}
\pf
Let $\fa=\fp_1^{e_1}\cdots\fp_r^{e_r}$ (9.4; if $\fa=A$ then $A/\fa=0$ and there is nothing to prove). Ideals of $A/\fa$ are of the form $\fb/\fa$ with $\fb\supseteq\fa$ an ideal of $A$. If $\fb\supseteq\fa$ then $\fb\ne0$ and $\fb^{-1}\fa\subseteq\fb^{-1}\fb=A$, so $\fa=\fb\cdot(\fb^{-1}\fa)$ is a multiple of $\fb$; by unique factorization $\fb=\fp_1^{f_1}\cdots\fp_r^{f_r}$ with $0\le f_i\le e_i$ (only primes dividing $\fa$ occur in $\fb$).

\emph{Choice of a generator.} For each $i$, $\fp_i^{f_i}\ne\fp_i^{f_i+1}$ (unique factorization), so pick $x_i\in\fp_i^{f_i}-\fp_i^{f_i+1}$. The ideals $\fp_i^{f_i+1}$ are pairwise coprime, so by the Chinese remainder theorem (1.10) there is $x\in A$ with $x\equiv x_i\pmod{\fp_i^{f_i+1}}$ for all $i$. Then $x\in\fp_i^{f_i}$ and $x\notin\fp_i^{f_i+1}$ (since $x-x_i\in\fp_i^{f_i+1}$ and $x_i\in\fp_i^{f_i}-\fp_i^{f_i+1}$).

\emph{Claim: $\fb=(x)+\fa$.} Both are ideals, so we compare localizations at every maximal ideal $\fm$ (Conventions). Recall that $\fp^{k+1}A_\fp\cap A=\fp^{k+1}$ for a non-zero prime $\fp$ (the ideal $\fp^{k+1}$ is $\fp$-primary, 4.2, and 4.8(ii)).
\begin{itemize}[leftmargin=*]
\item $\fm=\fp_i$: in the DVR $A_{\fp_i}$ the element $x$ lies in $\fp_i^{f_i}A_{\fp_i}$ but not in $\fp_i^{f_i+1}A_{\fp_i}$ (otherwise $x\in\fp_i^{f_i+1}A_{\fp_i}\cap A=\fp_i^{f_i+1}$), so $(x)_{\fm}=\fp_i^{f_i}A_{\fm}$. Moreover $\fa_{\fm}=\fp_i^{e_i}A_\fm$ ($\fp_j\not\subseteq\fm$ for $j\ne i$, so $\fp_j$ becomes the unit ideal), and $\fp_i^{e_i}A_\fm\subseteq\fp_i^{f_i}A_\fm$. So $((x)+\fa)_\fm=\fp_i^{f_i}A_\fm=\fb_\fm$.
\item $\fm\notin\{\fp_1,\dots,\fp_r\}$: then $\fa_\fm=A_\fm$ and $\fb_\fm=A_\fm$ (as $\fb\supseteq\fa$), so both sides equal $A_\fm$.
\end{itemize}
Thus $\fb=(x)+\fa$, i.e.\ $\fb/\fa$ is generated by the class of $x$ in $A/\fa$. Every ideal of $A/\fa$ is principal.

\emph{Every ideal is $2$-generated.} Let $\fb$ be an ideal of $A$. If $\fb=0$ it is principal. Otherwise pick $0\ne a\in\fb$ and apply the above to $\fa=(a)\ne0$: the ideal $\fb/(a)$ of $A/(a)$ is principal, generated by the class of some $x\in\fb$; hence $\fb=(a,x)$. \hfill$\square$

%---------------------------------------------------------------- 8
\exercise{8}{Let $\fa,\fb,\fc$ be three ideals in a Dedekind domain. Prove that
\begin{align*}
\fa\cap(\fb+\fc)&=(\fa\cap\fb)+(\fa\cap\fc),\\
\fa+(\fb\cap\fc)&=(\fa+\fb)\cap(\fa+\fc).
\end{align*}
[Localize.]}
\pf
In each identity the right-hand side is contained in the left-hand side for the first, and the left-hand side is contained in the right-hand side for the second, in any ring: $(\fa\cap\fb)+(\fa\cap\fc)\subseteq\fa\cap(\fb+\fc)$, and $\fa+(\fb\cap\fc)\subseteq(\fa+\fb)\cap(\fa+\fc)$. So it suffices to compare localizations at every maximal ideal $\fm$ (Conventions). Localization commutes with finite sums and intersections (3.4), so we must check both identities for the three ideals $\fa_\fm,\fb_\fm,\fc_\fm$ of the DVR $A_\fm$.

In a DVR with maximal ideal $(\pi)$ the ideals are $(\pi^k)$ ($k\ge0$) and $0$, so any two ideals are comparable under inclusion. By symmetry in $\fb,\fc$ we may assume $\fb\subseteq\fc$. Then $\fb+\fc=\fc$ and $\fb\cap\fc=\fb$. First identity: $\fa\cap(\fb+\fc)=\fa\cap\fc$, and $\fa\cap\fb\subseteq\fa\cap\fc$ gives $(\fa\cap\fb)+(\fa\cap\fc)=\fa\cap\fc$. Second identity: $\fa+(\fb\cap\fc)=\fa+\fb$, and $\fa+\fb\subseteq\fa+\fc$ gives $(\fa+\fb)\cap(\fa+\fc)=\fa+\fb$. So both identities hold locally, hence globally. \hfill$\square$

%---------------------------------------------------------------- 9
\exercise{9}{(\emph{Chinese Remainder Theorem}.) Let $\fa_1,\dots,\fa_n$ be ideals and let $x_1,\dots,x_n$ be elements in a Dedekind domain $A$. Then the system of congruences $x\equiv x_i\pmod{\fa_i}$ $(1\le i\le n)$ has a solution $x$ in $A$ $\iff$ $x_i\equiv x_j\pmod{\fa_i+\fa_j}$ whenever $i\ne j$.

[This is equivalent to saying that the sequence of $A$-modules
\[
A\xrightarrow{\ \phi\ }\bigoplus_{i=1}^nA/\fa_i\xrightarrow{\ \psi\ }\bigoplus_{i<j}A/(\fa_i+\fa_j)
\]
is exact, where $\phi$ and $\psi$ are defined as follows: $\phi(x)=(x+\fa_1,\dots,x+\fa_n)$; $\psi(x_1+\fa_1,\dots,x_n+\fa_n)$ has $(i,j)$-component $x_i-x_j+\fa_i+\fa_j$. To show that this sequence is exact it is enough to show that it is exact when localized at any $\fp\ne0$: in other words we may assume that $A$ is a discrete valuation ring, and then it is easy.]}
\pf
\emph{The reformulation.} A tuple $(x_i+\fa_i)_i$ lies in $\Ker\psi$ iff $x_i-x_j\in\fa_i+\fa_j$ for all $i<j$, and lies in $\operatorname{Im}\phi$ iff there is $x\in A$ with $x\equiv x_i\pmod{\fa_i}$ for all $i$. So the statement says exactly that $\Ker\psi=\operatorname{Im}\phi$. Always $\psi\phi=0$: the $(i,j)$-component of $\psi\phi(x)$ is $x-x$ modulo $\fa_i+\fa_j$. (This is also the easy direction of the equivalence: a solution $x$ gives $x_i-x_j\equiv x-x=0$.) Put $E=\Ker\psi/\operatorname{Im}\phi$; we must show $E=0$.

\emph{Localization.} Localization is exact (3.3), so for every maximal ideal $\fm$ the localized sequence is exact in the middle iff $E_\fm=0$ (because $(\Ker\psi)_\fm=\Ker\psi_\fm$ and $(\operatorname{Im}\phi)_\fm=\operatorname{Im}\phi_\fm$). Using $(A/\fa)_\fm\cong A_\fm/\fa_\fm$ and $(\fa_i+\fa_j)_\fm=(\fa_i)_\fm+(\fa_j)_\fm$ (3.4), the localized sequence is the same sequence for the ring $A_\fm$ and the ideals $(\fa_i)_\fm$. If $E_\fm=0$ for all maximal $\fm$, then $E=0$ (3.8). Every maximal ideal $\fm$ is non-zero, so $A_\fm$ is a DVR.

\emph{The DVR case.} Let $A$ be a DVR with maximal ideal $(\pi)$. Each $\fa_i$ is $0$ or $(\pi^{k_i})$; write $\fa_i=(\pi^{k_i})$ with $k_i\in\{0,1,2,\dots\}\cup\{\infty\}$ ($\pi^\infty:=0$). Then $\fa_i+\fa_j=(\pi^{\min(k_i,k_j)})$. Suppose $x_i-x_j\in\fa_i+\fa_j$ for all $i,j$. Renumber so that $k_1\le k_2\le\dots\le k_n$ and take $x=x_n$. For $i\le n$, $x_n-x_i\in\fa_i+\fa_n=(\pi^{\min(k_i,k_n)})=(\pi^{k_i})=\fa_i$ (for $i=n$ this is $0$). So $x\equiv x_i\pmod{\fa_i}$ for all $i$: a solution exists.

Hence $E_\fm=0$ for every maximal $\fm$, so $E=0$ and the system is solvable iff the compatibility conditions hold. \hfill$\square$

\end{document}
